\section{Superficies parametrizadas}
\label{sec:superficies}

\begin{definicion}[Superficie parametrizada]
  \label{def:superficie}
  Una \textbf{superficie parametrizada} es un mapeo
  \begin{equation}
    \label{eq:superficie}
    \vb{M}:\Om\subset\R^2\longrightarrow\R^m,
    \qquad m=2,3,\dots
  \end{equation}
  En el caso $m=3$,
  \begin{equation}
    \label{eq:superficie_3d}
    \vb{M}(\xii,\et)=
    \begin{pmatrix}
      M_1(\xii,\et)\\
      M_2(\xii,\et)\\
      M_3(\xii,\et)
    \end{pmatrix}.
  \end{equation}
\end{definicion}

\begin{ejemplo}[Esfera]
  \label{ej:esfera_superficie}
  \begin{equation}
    \label{eq:esfera_superficie}
    \vb{M}(\theta,\vph)=
    \begin{pmatrix}
      R\sin\theta\cos\vph\\
      R\sin\theta\sin\vph\\
      R\cos\theta
    \end{pmatrix}.
  \end{equation}
\end{ejemplo}

\begin{ejemplo}[Cilindro]
  \label{ej:cilindro}
  \begin{equation}
    \label{eq:cilindro}
    \vb{M}(\vph,z)=
    \begin{pmatrix}
      R\cos\vph\\
      R\sin\vph\\
      z
    \end{pmatrix},
    \quad \vph\in[0,2\pi),\ z\in[-h,h].
  \end{equation}
\end{ejemplo}

\subsection{Diferencial de superficie}
\label{subsec:diferencial_superficie}

Sea $\vb{M}:\Om\subset\R^2\rightarrow\R^3$, de modo que
$\vb{M}=\vb{M}(\xii,\et)$. Si fijamos $\et=\et_0$ y tomamos el mapeo como
funci\'on de $\xii$, obtenemos
\begin{equation}
  \label{eq:curva_xi}
  \vb{M}_{\xii}=\vb{M}(\xii,\et_0),
\end{equation}
que es una curva sobre la superficie. An\'alogamente,
\begin{equation}
  \label{eq:curva_eta}
  \vb{M}_{\et}=\vb{M}(\xii_0,\et)
\end{equation}
es otra curva sobre $M$ en la direcci\'on $\et$. Los vectores tangentes a esas
curvas son
\begin{equation}
  \label{eq:tangentes_superficie}
  \vb{T}_{\xii}=\pdv{}{\xii}\vb{M}(\xii,\et_0),
  \qquad
  \vb{T}_{\et}=\pdv{}{\et}\vb{M}(\xii_0,\et),
\end{equation}
y la normal al plano tangente es su producto cu\~na,
\begin{equation}
  \label{eq:normal_superficie}
  \vb{N}=\eval{\pdv{}{\xii}\vb{M}\wedge\pdv{}{\et}\vb{M}}
    _{(\xii_0,\et_0)},
\end{equation}
de modo que el plano tangente en $\vb{M}_0$ es
$\vb{N}\cdot\qty{\vb{r}-\vb{M}_0}=0$.

La consecuencia importante es que las diferenciales de esas curvas, que viven
en el plano tangente, aproximan localmente a la superficie. En efecto, para
$\vb{M}(\xii(t),\et_0)$ y $\vb{M}(\xii_0,\et(t))$ la regla de la cadena da
\begin{equation}
  \label{eq:derivadas_curvas_superficie}
  \dv{}{t}\vb{M}_{\xii}=\pdv{}{\xii}\vb{M}\dv{\xii}{t},
  \qquad
  \dv{}{t}\vb{M}_{\et}=\pdv{}{\et}\vb{M}\dv{\et}{t},
\end{equation}
y por lo tanto
\begin{equation}
  \label{eq:diferenciales_curvas}
  \dd\vb{M}_{\xii}=\pdv{}{\xii}\vb{M}\dd\xii,
  \qquad
  \dd\vb{M}_{\et}=\pdv{}{\et}\vb{M}\dd\et .
\end{equation}

\begin{definicion}[Diferencial de superficie]
  \label{def:dS}
  \begin{equation}
    \label{eq:dS}
    \dd\vb{S}=\pdv{}{\xii}\vb{M}\wedge\pdv{}{\et}\vb{M}
      \dd\xii\dd\et ,
  \end{equation}
  y el \textbf{elemento de \'area} es su norma,
  \begin{equation}
    \label{eq:dS_norma}
    \dd S=\norm{\dd\vb{S}}
      =\norm{\pdv{}{\xii}\vb{M}\wedge\pdv{}{\et}\vb{M}}
       \abs{\dd\xii\dd\et} .
  \end{equation}
\end{definicion}

\subsection{Ejemplo: el elipsoide de revoluci\'on}
\label{subsec:elipsoide}

Sea $\vb{M}:\R^2\rightarrow\R^3$ con
\begin{equation}
  \label{eq:elipsoide_param}
  \vb{M}(\theta,\vph)=
  \begin{pmatrix}
    a\sin\theta\cos\vph\\
    a\sin\theta\sin\vph\\
    b\cos\theta
  \end{pmatrix}.
\end{equation}
Entonces
\begin{equation}
  \label{eq:dS_elipsoide}
  \small
  \begin{aligned}
    \dd\vb{S}
      &=\pdv{\vb{M}}{\theta}\wedge\pdv{\vb{M}}{\vph}
        \dd\theta\dd\vph\\
      &=
      \begin{pmatrix}
        a\cos\theta\cos\vph\\
        a\cos\theta\sin\vph\\
        -b\sin\theta
      \end{pmatrix}
      \wedge
      \begin{pmatrix}
        -a\sin\theta\sin\vph\\
        a\sin\theta\cos\vph\\
        0
      \end{pmatrix}
      \dd\theta\dd\vph\\
      &=
      \begin{pmatrix}
        ab\sin^2\theta\cos\vph\\
        ab\sin^2\theta\sin\vph\\
        a^2\sin\theta\cos\theta
      \end{pmatrix}
      \dd\theta\dd\vph\\
      &=a\sin\theta
      \begin{pmatrix}
        b\sin\theta\cos\vph\\
        b\sin\theta\sin\vph\\
        a\cos\theta
      \end{pmatrix}
      \dd\theta\dd\vph .
  \end{aligned}
\end{equation}

En particular, si $b=a$ recuperamos la esfera:
\begin{equation}
  \label{eq:dS_esfera}
  \dd\vb{S}=\vh{r}\,a^2\sin\theta\dd\theta\dd\vph,
  \qquad
  \vh{r}=
  \begin{pmatrix}
    \sin\theta\cos\vph\\
    \sin\theta\sin\vph\\
    \cos\theta
  \end{pmatrix}.
\end{equation}

\subsubsection{Integraci\'on del vector de \'area}
\label{subsubsec:integral_dS}

Integrando \eqref{eq:dS_elipsoide} sobre toda la superficie, las integrales en
$\vph$ de $\cos\vph$ y $\sin\vph$ se anulan y queda
\begin{equation}
  \label{eq:integral_dS}
  \begin{aligned}
    \iint\dd\vb{S}
      &=\int_{0}^{2\pi}\!\!\int_{0}^{\pi}
        \begin{pmatrix}
          b\sin\theta\cos\vph\\
          b\sin\theta\sin\vph\\
          a\cos\theta
        \end{pmatrix}
        a\sin\theta\dd\theta\dd\vph\\
      &=2\pi a^2\vh{z}\int_{0}^{\pi}\sin\theta\cos\theta\dd\theta
       =\vb{0},
  \end{aligned}
\end{equation}
como debe ser para toda superficie cerrada.

\subsubsection{Integral de \'area}
\label{subsubsec:integral_area}

\begin{equation}
  \label{eq:area_elipsoide}
  \begin{aligned}
    A&=\iint\norm{\dd\vb{S}}=\iint\dd S\\
     &=\int_{0}^{2\pi}\!\!\int_{0}^{\pi}
       a\sqrt{b^2\sin^2\theta+a^2\cos^2\theta}\,
       \sin\theta\dd\theta\dd\vph\\
     &=2\pi a\int_{0}^{\pi}
       \sqrt{b^2+(a^2-b^2)\cos^2\theta}\,\sin\theta\dd\theta .
  \end{aligned}
\end{equation}
En particular, si $a=b$,
\begin{equation}
  \label{eq:area_esfera}
  A=2\pi a^2\int_{0}^{\pi}\sin\theta\dd\theta=4\pi a^2,
\end{equation}
el \'area de la esfera.

\subsection{Diferencial de volumen}
\label{subsec:diferencial_volumen}

Supongamos ahora un mapeo $\vb{M}:\R^3\rightarrow\R^3$ con par\'ametros
$u,v,w$. Los vectores tangentes a las curvas coordenadas son
\begin{equation}
  \label{eq:tangentes_volumen}
  \vb{M}_u=\pdv{}{u}\vb{M},
  \qquad
  \vb{M}_v=\pdv{}{v}\vb{M},
  \qquad
  \vb{M}_w=\pdv{}{w}\vb{M} .
\end{equation}
Como el volumen de la caja que generan tres vectores es
$(\vb{a}\wedge\vb{b})\cdot\vb{c}$, el elemento de volumen es
\begin{equation}
  \label{eq:dV}
  \dd V=\qty{\vb{M}_u\wedge\vb{M}_v}\cdot\vb{M}_w
    \dd u\dd v\dd w .
\end{equation}

\begin{ejemplo}[Volumen del elipsoide]
  \label{ej:volumen_elipsoide}
  Tomemos el mapeo de la caja al elipsoide
  $\vb{M}=\vb{M}(\xii,\theta,\vph)$ con
  \begin{equation}
    \label{eq:mapeo_elipsoide}
    \vb{M}=
    \begin{pmatrix}
      a\xii\sin\theta\cos\vph\\
      b\xii\sin\theta\sin\vph\\
      c\xii\cos\theta
    \end{pmatrix},
    \qquad a<b<c .
  \end{equation}
  Un c\'alculo directo del triple producto da
  \begin{equation}
    \label{eq:dV_elipsoide}
    \dd V=\qty{\pdv{\vb{M}}{\xii}\wedge\pdv{\vb{M}}{\theta}}
      \cdot\pdv{\vb{M}}{\vph}\dd\xii\dd\theta\dd\vph
      =abc\,\xii^2\sin\theta\dd\xii\dd\theta\dd\vph ,
  \end{equation}
  de modo que
  \begin{equation}
    \label{eq:volumen_elipsoide}
    \begin{aligned}
      \mathrm{Vol}&=abc\int_{0}^{2\pi}\!\!\int_{0}^{\pi}\!\!\int_{0}^{1}
        \xii^2\sin\theta\dd\xii\dd\theta\dd\vph\\
        &=abc\,(2\pi)(2)\qty{\tfrac{1}{3}}=\frac{4\pi}{3}abc .
    \end{aligned}
  \end{equation}
  Si $a=b=c$ se recupera $\mathrm{Vol}=\tfrac{4}{3}\pi a^3$.
\end{ejemplo}

\begin{observacion}[Sistema positivo]
  Al construir $\dd\vb{S}$ o $\dd V$ hay que \emph{preservar el orden de las
  variables}: el signo del producto cu\~na y del triple producto depende de
  \'el, y con ello la orientaci\'on de la normal o el signo del volumen.
\end{observacion}
